\question[4]
Gegeben ist der folgende Python-Code:

\begin{lstlisting}
bonus = 5

def punkte_berechnen(punkte):
    gesamt = punkte + bonus
    return gesamt

ergebnis = punkte_berechnen(20)
print(ergebnis)
print(bonus)
print(gesamt)
\end{lstlisting}

\begin{parts}
\part Gib alle globalen und alle lokalen Variablen an.
\part Gib an, was die ersten beiden \texttt{print}-Anweisungen ausgeben.
\part Die letzte Zeile führt zu einem Fehler. Erläutere mit dem Begriff Gültigkeitsbereich, warum.
\end{parts}

\begin{solutionbox}{6cm}
\begin{parts}
\part Global: \texttt{bonus}, \texttt{ergebnis}. Lokal (in der Funktion): \texttt{punkte}, \texttt{gesamt}.
\part \texttt{25} und \texttt{5}
\part \texttt{gesamt} ist eine lokale Variable der Funktion. Ihr Gültigkeitsbereich ist nur die Funktion, und sie existiert
nur, solange die Funktion ausgeführt wird. Außerhalb ist sie nicht bekannt, deshalb gibt es einen \texttt{NameError}.
\end{parts}
\end{solutionbox}
